\documentclass[11pt,a4paper]{article}

\usepackage{fontspec}
\usepackage{polyglossia}
\setdefaultlanguage{turkish}
\setmainfont{Carlito}
\setmonofont[Scale=0.86]{DejaVu Sans Mono}

\usepackage[a4paper,margin=2.4cm,headsep=0.7cm]{geometry}
\usepackage{graphicx}
\usepackage{longtable}
\usepackage{booktabs}
\usepackage{array}
\usepackage{enumitem}
\usepackage{fancyhdr}
\usepackage{xcolor}
\usepackage{titlesec}
\usepackage{caption}
\usepackage{float}

\definecolor{koyu}{RGB}{30,45,70}
\definecolor{gri}{RGB}{110,110,110}

\titleformat{\section}{\large\bfseries\color{koyu}}{\thesection.}{0.6em}{}
\titlespacing{\section}{0pt}{1.4em}{0.7em}

\captionsetup{font=small,labelfont={bf,color=koyu},skip=6pt}

\pagestyle{fancy}
\fancyhf{}
\lhead{\small\color{gri}Sistem Algoritması Tasarımı}
\rhead{\small\color{gri}\thepage}
\renewcommand{\headrulewidth}{0.4pt}

\setlength{\parskip}{0.55em}
\setlength{\parindent}{0pt}
\setlist{itemsep=0.25em,topsep=0.3em}

\newcommand{\ey}[1]{\texttt{#1}}

\begin{document}

\begin{center}
{\Large\bfseries\color{koyu} Sistem Algoritması Tasarımı}\\[0.4em]
{\large Robot Kol ve Otonom Araç}\\[1.1em]
{\small TEKNOFEST 2026 Mesleki Yetenek Yarışması\\
Akıllı Fabrika Sistemleri Programlama Kategorisi, Yarı Final Görev 1}\\[0.6em]
{\small\color{gri} Takım: andromeda}
\end{center}

\vspace{1.2em}

Bu doküman robot kol ile otonom aracı birlikte çalıştıran algoritmayı anlatıyor.
İçinde sistemin genel işleyişi, dört akış şeması, tam sözde kod, tasarım
kararlarının gerekçeleri ve dayandığımız varsayımlar var.

Tasarımda yalnızca bize verilen hazır eylemleri kullandık. Bunların dışında
hiçbir donanım çağrısı, sensör erişimi ya da kütüphane fonksiyonu varsaymadık.

\vspace{0.8em}

\section{Sistemin genel işleyişi}

Sistemde üç ayrı denetleyici var ve bunlar iki farklı kanaldan haberleşiyor.
PLC hattı yönetiyor ve robot kola tek bir dijital başla sinyali gönderiyor.
Sinyali K10 rölesinin kuru kontağı taşıyor, Arduino'nun A5 ucunu şaseye
çekiyor. Robot kol küpün
rengini görüntü işlemeyle belirliyor, küpü araca yüklüyor ve rengi MQTT
üzerinden \ey{arac/yuk} konusuna yayınlıyor. Otonom araç da bu rengi dinliyor,
küpün üzerinde olduğunu MZ80 ile doğruluyor, sonra pist görevlerini tamamlayıp
renge uygun park alanına giriyor.

İki sürecin ömrü farklı ve bu ayrım tasarımın temeli. Robot kol süreci sonsuz
döngü, konveyöre küp geldikçe çevrim tekrarlanıyor. Otonom araç süreci ise tek
görevlik, park etmesiyle birlikte sona eriyor.

\begin{figure}[H]
\centering
\includegraphics[width=0.97\textwidth]{sema_sistem.pdf}
\caption{Sistem etkileşim diyagramı}
\end{figure}

\clearpage

\section{Kullandığımız hazır eylemler}

Aşağıdaki iki tabloda bize verilen hazır eylemlerin tamamı ve her birini bu
tasarımda nerede kullandığımız var. Tasarım bu kümenin dışına çıkmıyor.

\subsection*{Robot kol}

\begin{longtable}{@{}p{3.4cm}p{4.1cm}p{7.2cm}@{}}
\toprule
\textbf{Eylem} & \textbf{Ne yapar} & \textbf{Nerede kullandık} \\
\midrule
\endhead
\ey{home()} & Tüm eksenleri referans konumuna götürür. &
Sadece başlangıçta bir kez. Güç açılışında kolun nerede durduğunu bilmiyoruz. \\
\addlinespace
\ey{git("GORME")} & Öğretilmiş poza gider. &
Bekleme pozu, kamera buradan konveyör çıkışını görüyor. Yüklemeden sonra da
güvenli poz olarak aynı yere dönüyoruz. \\
\addlinespace
\ey{git("AL")} & & Küpü kavrama pozu. \\
\addlinespace
\ey{git("YUKLE")} & & Küpü araca bırakma pozu. \\
\addlinespace
\ey{tut()} & Tutucuyu kapatır. & AL pozunda, küpü kavramak için. \\
\addlinespace
\ey{birak()} & Tutucuyu açar. &
YUKLE pozunda küpü bırakırken. Ayrıca başlangıçta tutucuyu boşaltmak için. \\
\addlinespace
\ey{renk\_oku()} & RED, GREEN, BLUE veya UNKNOWN döndürür. &
GORME pozunda üç kez çağırıp çoğunluk oyu alıyoruz (Şekil 4). \\
\addlinespace
\ey{basla\_geldi()} & PLC başla sinyali, True veya False. &
Çevrim başında beklemek için, bir de çevrim sonunda sinyalin düşmesini
beklemek için. \\
\addlinespace
\ey{araca\_gonder(renk)} & Rengi MQTT ile araca bildirir. &
Çevrimin en son adımı, kol güvenli poza çekildikten sonra. \\
\bottomrule
\end{longtable}

\subsection*{Otonom araç}

\begin{longtable}{@{}p{3.4cm}p{4.1cm}p{7.2cm}@{}}
\toprule
\textbf{Eylem} & \textbf{Ne yapar} & \textbf{Nerede kullandık} \\
\midrule
\endhead
\ey{yuk\_rengi\_bekle()} & MQTT'den renk gelene kadar bekler. &
Birinci şart. Bloklayıcı olduğu için süreç burada duruyor. \\
\addlinespace
\ey{yuk\_algilandi()} & MZ80 küpü gördüyse True. &
İkinci şart. Sağlanana kadar yokluyoruz, araç hareket etmiyor. \\
\addlinespace
\ey{ilerle()} & Aracı ilerletir. &
Yalnızca iki şart birlikte sağlandıktan sonra, bir de her duruşun ardından. \\
\addlinespace
\ey{dur()} & Aracı durdurur. &
Başlangıçta, kırmızı ışıkta, yaya ve tümsek görevlerinde, park sonunda. \\
\addlinespace
\ey{bekle(sn)} & Verilen süre kadar bekletir. &
Yoklama aralığı, tabelaya yaklaşma süresi ve görev süreleri. \\
\addlinespace
\ey{isik\_oku()} & KIRMIZI, YESIL veya YOK döndürür. &
Trafik ışığı alt yordamında (Şekil 5). \\
\addlinespace
\ey{tabela\_oku()} & YAYA, TUMSEK, PARK veya YOK döndürür. &
Görev döngüsünün her turunda, kilitleme mantığıyla (Şekil 3). \\
\addlinespace
\ey{park\_et(renk)} & Verilen renkteki park alanına park eder. &
PARK tabelası işlendiğinde, görevin son adımı. \\
\bottomrule
\end{longtable}

\subsection*{Sabitler}

\begin{longtable}{@{}p{4.2cm}p{1.5cm}p{9cm}@{}}
\toprule
\textbf{Sabit} & \textbf{Değer} & \textbf{Neden bu değer} \\
\midrule
\endhead
\ey{YAKLASMA\_SN} & 1,5 s &
Tabela kamerada belli bir mesafeden görünüyor. Sabit süre ilerlemek görevi her
seferinde aynı noktada uygulamamızı sağlıyor. Saha denemesiyle ayarlanacak. \\
\addlinespace
\ey{YAYA\_BEKLEME\_SN} & 3,0 s &
Yaya geçidindeki tam duruşun dışarıdan net görülebilmesi için. \\
\addlinespace
\ey{TUMSEK\_BEKLEME\_SN} & 1,0 s &
Hız kesmenin eldeki eylemlerle en yakın karşılığı. Varsayım V-3'e bakınız. \\
\addlinespace
\ey{RENK\_DENEME\_MAX} & 5 &
Renk okumada tekrar üst sınırı, sonsuz döngüye girmemek için. \\
\addlinespace
\ey{YOKLAMA\_SN} & 0,1 s &
Beklerken işlemciyi meşgul etmemek için. 100 ms'lik tepki gecikmesi bu hat
hızında fark etmiyor. \\
\bottomrule
\end{longtable}

\clearpage

\section{Robot kol akış şeması}

\begin{figure}[H]
\centering
\includegraphics[height=0.82\textheight,keepaspectratio]{sema_robot.pdf}
\caption{Robot kol akış şeması}
\end{figure}

\clearpage

\section{Otonom araç akış şeması}

\begin{figure}[H]
\centering
\includegraphics[height=0.82\textheight,keepaspectratio]{sema_arac.pdf}
\caption{Otonom araç akış şeması}
\end{figure}

\clearpage

\section{Alt yordamlar}

\begin{figure}[H]
\centering
\includegraphics[width=0.6\textwidth]{sema_renk.pdf}
\caption{Alt yordam: \ey{RENK\_OKU\_GUVENLI()}}
\end{figure}

\vspace{1.2em}

\begin{figure}[H]
\centering
\includegraphics[width=0.6\textwidth]{sema_isik.pdf}
\caption{Alt yordam: \ey{TRAFIK\_ISIGI\_KONTROL()}}
\end{figure}

\clearpage

\section{Sözde kod}

Aşağıdaki sözde kod Şekil 2 ve Şekil 3'teki akış şemalarının birebir karşılığı.
İki süreç fiziksel olarak ayrı denetleyicilerde eşzamanlı koşuyor: kol tarafında
Arduino ile Jetson, araç tarafında Jetson ile ESP32. Aralarındaki tek bağ MQTT
mesajı.

\subsection*{Sabitler}

\begin{verbatim}
SABITLER
    YAKLASMA_SN        <- 1.5   # tabela goruldukten sonra yaklasma suresi
    YAYA_BEKLEME_SN    <- 3.0   # yaya gecidinde tam durus
    TUMSEK_BEKLEME_SN  <- 1.0   # tumsekte hiz kesme
    RENK_DENEME_MAX    <- 5
    YOKLAMA_SN         <- 0.1
    GECERLI_RENKLER    <- { "RED", "GREEN", "BLUE" }
\end{verbatim}

\subsection*{Süreç A: robot kol, sonsuz döngü}

\begin{verbatim}
YORDAM ROBOT_KOL_ANA()

    # ---- Baslangic. Dongunun disinda, yalnizca bir kez ---------------
    home()                    # konum bilinmiyor, referansa git
    birak()                   # tutucuda kup kalmis olabilir
    git("GORME")              # kameranin konveyoru gordugu bekleme pozu

    DONGU SONSUZ                                                   # (A)

        # ---- 1) PLC'nin basla sinyalini bekle --------------------
        # Sinyal, kup konveyor cikisindayken VE arac yukleme
        # noktasinda dururken geliyor. Bu sarti PLC garanti ediyor,
        # kol ayrica kontrol etmiyor.
        WHILE NOT basla_geldi()
            bekle(YOKLAMA_SN)
        END WHILE

        # ---- 2) Rengi guvenli bicimde oku ------------------------
        renk <- RENK_OKU_GUVENLI()

        EGER renk GECERLI_RENKLER icinde DEGILSE
            # Yanlis renk yanlis park alani demek. Tahmin etmek
            # yerine cevrimi kesiyoruz.
            HATA_BILDIR("renk belirlenemedi")
            # kup ALINMAZ, araca hicbir sey GONDERILMEZ

        DEGILSE
            # ---- 3) Kupu al ve araca yukle ----------------------
            git("AL")
            tut()
            git("YUKLE")
            birak()

            # ---- 4) Kolu guvenli poza cek, ANCAK SONRA bildir --
            # Sira kritik: arac, renk bildirimini alir almaz
            # hareket edebiliyor.
            git("GORME")
            araca_gonder(renk)
        END EGER

        # ---- 5) Ayni sinyalle ikinci cevrimi onle ---------------
        # PLC sinyali SEVIYE olarak suruyor. Dusmesini beklemezsek
        # bos bir cevrim daha baslatiyoruz.
        WHILE basla_geldi()
            bekle(0.05)
        END WHILE

    END DONGU                                     # -> (A) dongu basi

END YORDAM
\end{verbatim}

\subsection*{Alt yordam: RENK\_OKU\_GUVENLI}

\begin{verbatim}
ALT YORDAM RENK_OKU_GUVENLI() -> renk
    deneme <- 0
    WHILE deneme < RENK_DENEME_MAX
        o1 <- renk_oku()
        o2 <- renk_oku()
        o3 <- renk_oku()
        c  <- COGUNLUK(o1, o2, o3)   # ucunden en az ikisinde ayni
                                     # olan deger, yoksa YOK doner
        EGER c, YOK ve "UNKNOWN" degilse ISE
            DONDUR c
        END EGER
        deneme <- deneme + 1
    END WHILE
    DONDUR "UNKNOWN"           # 5 denemede de kararsiz, hata cagrisi
END ALT YORDAM
\end{verbatim}

\clearpage

\subsection*{Süreç B: otonom araç, tek görevlik}

\begin{verbatim}
YORDAM OTONOM_ARAC_ANA()

    dur()                         # baslangicta kesinlikle hareketsiz

    # ---- Birinci sart: MQTT'den renk bildirimi, bloklayici -------
    renk <- yuk_rengi_bekle()

    # ---- Ikinci sart: MZ80 kupu fiziksel olarak algilamali -------
    WHILE NOT yuk_algilandi()
        dur()
        bekle(YOKLAMA_SN)
    END WHILE

    # Buraya ancak IKI SART DA saglandiginda geliniyor.
    ilerle()
    son_tabela <- "YOK"

    DONGU SONSUZ                                                # (C)

        # ---- Trafik isigi her turda kontrol ediliyor ---------
        TRAFIK_ISIGI_KONTROL()

        t <- tabela_oku()

        # ---- Kilit acma: tabela gorus alanindan cikti --------
        EGER t = "YOK" ISE
            son_tabela <- "YOK"
            DEVAM
        END EGER

        # ---- Kilit: ayni tabela ikinci kez islenmiyor --------
        # Bir tabela kamerada onlarca kare gorunuyor. Kilit
        # olmadan arac ayni gorevi defalarca uyguluyor.
        EGER t = son_tabela ISE
            DEVAM
        END EGER

        son_tabela <- t                                  # kilitle
        bekle(YAKLASMA_SN)        # tabela uzaktan gorunur, yaklas

        SEC t
            DURUM "PARK":
                dur()
                park_et(renk)     # renk birinci sarttan geldi
                dur()
                CIK DONGU              # arac TEK GOREVLIKTIR

            DURUM "YAYA":
                dur()
                bekle(YAYA_BEKLEME_SN)      # tam durus
                ilerle()

            DURUM "TUMSEK":
                dur()
                bekle(TUMSEK_BEKLEME_SN)    # hiz kesme, bkz. V-3
                ilerle()
        END SEC

    END DONGU                              # -> (C) gorev dongusu

    # Gorev tamamlandi, surec burada sona eriyor.

END YORDAM
\end{verbatim}

\subsection*{Alt yordam: TRAFIK\_ISIGI\_KONTROL}

\begin{verbatim}
ALT YORDAM TRAFIK_ISIGI_KONTROL()

    EGER isik_oku() "KIRMIZI" DEGILSE
        DONDUR             # isik yok veya yesil, ilerlemeye devam
    END EGER

    # Kirmizi gorundu.
    dur()
    WHILE isik_oku() "YESIL" DEGIL   # CIKIS SARTI YALNIZCA "YESIL"
        bekle(YOKLAMA_SN)            # "YOK" okumak GECME IZNI DEGIL
    END WHILE
    ilerle()

END ALT YORDAM
\end{verbatim}

\clearpage

\section{Tasarım kararları}

Görev tanımında cevaplanması istenen on iki karar noktası var. Aşağıda her
birini, verdiğimiz kararla birlikte neden başka türlü yapmadığımızı da
anlatarak yanıtladık.

\subsection*{Robot kol}

\textbf{Sistem ilk açıldığında, döngüye girmeden önce yapılması gereken bir şey
var mı?}

Var, üç adım: \ey{home()}, sonra \ey{birak()}, sonra \ey{git("GORME")}. Üçü de
döngünün dışında ve yalnızca bir kez çalışıyor.

\ey{home()} gerekiyor çünkü güç kesintisi ya da acil duruş sonrası kol herhangi
bir pozda kalmış olabilir. Öğretilmiş pozlara güvenli bir yörünge ancak bilinen
bir referans noktasından üretilebiliyor. Bu adımı atlarsak ilk hareket
öngörülemez bir yoldan geçiyor ve konveyöre ya da araca çarpma riski doğuyor.

\ey{birak()} koyduk çünkü önceki çevrim yükleme anında kesilmişse tutucuda küp
kalmış olabilir. Tutucu kapalıyken \ey{tut()} çağırmak yeni küpü kavramıyor ve
çevrim sessizce boşa dönüyor.

\ey{git("GORME")} ise renk okumasının bu pozda yapılacak olması yüzünden.
Böylece kol sinyali beklerken zaten doğru yerde duruyor. Sinyal geldikten sonra
poza gitmek çevrim süresine gereksiz hareket ekliyor, süre de finalde ayrı bir
puan kalemi.

\textbf{Başla sinyalini beklemeyi nasıl kurguladınız?}

Yoklamalı bekleme kullandık: \ey{WHILE NOT basla\_geldi(): bekle(0.1)}.

Sinyali PLC seviye olarak sürüyor, yani K10 rölesi robot çevrimi boyunca çekili
kalıyor. Dolayısıyla beklediğimiz şey sinyalin gelmesi değil, var olması.
\ey{bekle(0.1)} ile araya boşluk koymamız sıkı döngünün işlemciyi gereksiz
meşgul etmesini önlüyor. 100 ms'lik tepki gecikmesi bu hattın çevrim süresi
yanında önemsiz.

Beklemenin ikinci yarısı çevrimin sonunda: iş bittikten sonra sinyalin düşmesini
bekliyoruz. İkisi birlikte bir kenar yakalama davranışı üretiyor.

\textbf{Rengi ne zaman ve kaç kez okursunuz? UNKNOWN gelirse ne olur?}

Kol GORME pozundayken, yani küp konveyör çıkışında dururken ve kavramadan önce
okuyoruz. Kavradıktan sonra okumanın iki riski var: tutucu küpün bir yüzünü
kapatabiliyor, kolun gövdesi de küpün üzerine gölge düşürebiliyor. HSV eşikleri
gölgeye duyarlı olduğu için bu iki durum da okumayı bozuyor.

Üç kez okuyup çoğunluk oyu alıyoruz (Şekil 4). Tek okuma o anki parlamadan veya
otomatik pozlama dalgalanmasından etkilenebiliyor. Üç okumadan ikisinin aynı
çıkması rastgele gürültüyü büyük ölçüde eliyor, maliyeti de birkaç yüz
milisaniye.

UNKNOWN gelirse okumayı en fazla beş kez tekrarlıyoruz. Beş denemede de karar
verilemezse küpü almıyoruz ve araca hiçbir şey göndermiyoruz, çevrimi hata
bildirimiyle kesiyoruz.

Tahmin etmemeyi bilerek seçtik. Renk, aracın hangi park alanına gideceğini
belirliyor. Yanlış bir renk göndermek aracı yola çıkarıp görevi hatalı
tamamlatıyor. Bu, hiç başlatmamaktan daha kötü, çünkü hatayı ancak parkta fark
ediyoruz ve o noktada düzeltmek mümkün olmuyor. Ayrıca UNKNOWN çoğu zaman küpün
yanlış konumlandığını gösteriyor, bu da operatör müdahalesi isteyen fiziksel bir
sorun.

\textbf{Alma ve yükleme hangi sırayla yapılır?}

\ey{git("AL")}, \ey{tut()}, \ey{git("YUKLE")}, \ey{birak()}, \ey{git("GORME")}.

Her hareket komutu ile tutucu komutu arasında sıra bağımlılığı var. Kol hedefe
varmadan tutucu kapanırsa küp kavranmıyor, hedefe varmadan açılırsa küp yanlış
yere düşüyor. Bu yüzden \ey{tut()} yalnızca AL pozuna varıldıktan,
\ey{birak()} yalnızca YUKLE pozuna varıldıktan sonra çağrılıyor.

Sondaki \ey{git("GORME")} iki işi birden yapıyor: bir sonraki çevrimin bekleme
pozunu hazırlıyor ve kolu aracın hareket koridorundan çıkarıyor.

\textbf{Rengi araca ne zaman gönderirsiniz?}

Çevrimin en son adımında, küp bırakıldıktan ve kol GORME pozuna çekildikten
sonra. Tasarımdaki en kritik sıralama kararı bu.

Araç iki şartla hareket ediyor: MZ80'in küpü algılaması ve renk bildiriminin
gelmesi. Küp araca bırakıldığı anda MZ80 zaten aktif oluyor. Rengi o anda
gönderirsek ikinci şart da sağlanıyor ve araç hemen hareket ediyor. Oysa kol
henüz araç üzerinden çekilmemiş durumda. Sonuç: araç inen kolun altından
kayıyor, küp devriliyor veya kol ile araç çarpışıyor.

Bildirimi en sona almak bu çarpışmayı bir zamanlama ayarına değil, algoritmanın
yapısına bağlıyor. Bir faydası daha var: renk erken gönderilip kavrama
sonradan başarısız olsaydı araç boş yola çıkmış olurdu. Bu sıralama onu da
imkânsız kılıyor.

\textbf{Kol geri döndüğünde sinyal hâlâ aktifse veya değilse ne olur?}

Sinyal hâlâ aktifse bu yeni bir küp geldiği anlamına gelmiyor. PLC sinyali
seviye olarak sürüyor ve küp banttan kalkıp BP2 sensörü boşalana kadar çekili
kalabiliyor. Bu yüzden algoritma önce sinyalin düşmesini bekliyor:
\ey{WHILE basla\_geldi(): bekle(0.05)}. Bu adım olmasa aynı küp için ikinci bir
çevrim başlıyor, kol boş bir kavrama yapıyor ve araca ikinci kez renk
gönderiyor.

Sinyal düşmüşse doğrudan döngü başına dönüyoruz ve yeni sinyali bekliyoruz.
İkisi birlikte kenar yakalama davranışı üretiyor: düşen kenar, ardından
yükselen kenar. Seviye sinyaliyle çalışan her el sıkışmanın standart çözümü bu.

\textbf{Süreç nasıl tekrarlanır, nerede başlar nereye döner?}

Başlangıç bloğu döngünün dışında ve program ömrü boyunca bir kez çalışıyor.
Döngü Şekil 2'de A ile işaretli noktadan başlıyor, sinyalin düşmesi
beklendikten sonra yine A'ya dönüyor.

Hata dalı da aynı noktaya dönüyor ama doğrudan değil, önce sinyalin düşmesini
bekleyerek. Böylece hata sonrasında da aynı küp için sonsuz tekrar oluşmuyor.

Görev tanımı çevrimin konveyöre küp geldikçe sürekli tekrarlanmasını istediği
için döngünün çıkışı yok, süreç yalnızca güç kesildiğinde sona eriyor.

\subsection*{Otonom araç}

\textbf{Araç ne zaman harekete başlar? İki şarttan yalnızca biri sağlanırsa ne
olur?}

Araç yalnızca iki şart birlikte sağlandığında hareket ediyor. Birincisi
\ey{yuk\_rengi\_bekle()} fonksiyonunun MQTT'den bir renk döndürmesi, ikincisi
\ey{yuk\_algilandi()} değerinin True olması, yani MZ80'in küpü fiziksel olarak
görmesi.

Yalnızca renk gelirse küp araca yüklenmemiş demektir. Kavrama başarısız olmuş
ya da küp düşmüş olabilir. Araç boş yola çıkarsa görevi tamamlıyor ama yükü
taşımıyor, bu da görevin amacını ortadan kaldırıyor. O yüzden \ey{dur()}
konumunda bekliyor.

Yalnızca MZ80 algılarsa küp araçta ama rengini bilmiyoruz. Renk, park alanı
seçiminin tek girdisi, onsuz araç nereye gideceğini bilemiyor. Rastgele bir
parka gitmek beklemekten kötü olurdu, o yüzden bekliyor.

Algoritmadaki sıra bilinçli: önce bloklayıcı \ey{yuk\_rengi\_bekle()}, sonra
\ey{yuk\_algilandi()} yoklaması. Renk zaten kol tarafından en son gönderildiği
için, renk geldiğinde küp fiziksel olarak araçta olmuş oluyor. İkinci kontrol
bunun doğrulaması.

\textbf{Kırmızıda beklemeyi nasıl kurguladınız? Yeşile dönüşü nasıl
anlıyorsunuz?}

Ayrı bir alt yordamda ele aldık (Şekil 5). Mantık iki aşamalı. Okuma KIRMIZI
değilse hiçbir şey yapmıyoruz, araç ilerlemeye devam ediyor. Okuma KIRMIZI ise
\ey{dur()} çağırıp bekleme döngüsüne giriyoruz. Bu döngüden çıkış şartı yalnızca
YESIL okumak.

Buradaki kritik ayrım şu: YOK geçme izni değil. \ey{isik\_oku()} üç değer
döndürüyor ve YOK ışığın görülemediğini söylüyor, söndüğünü değil. Kırmızıda
beklerken araç titreşimi, güneş yansıması veya anlık bir algılama kaybı YOK
üretebiliyor. Çıkış şartını "KIRMIZI değil" diye yazsaydık araç kırmızı ışıkta
geçerdi. Bu kesin olarak yasak, o yüzden çıkış şartını olumlu biçimde YESIL
olarak yazdık.

Zaman aşımı bilerek koymadık. Işık arızalanırsa araç süresiz bekliyor.
Alternatif, yani belirli bir süre sonra geçmek, kural ihlali riskini yapısal
hâle getirirdi. Beklemek kural ihlalinden daha düşük maliyetli bir başarısızlık.

\textbf{Tabela görüldüğü an mı durulur, yaklaşınca mı? Aynı tabela iki kez
görülürse ne olur?}

Görüldüğü an durmuyoruz. Tabela kamerada belli bir mesafeden görünüyor, o anda
durursak yaya geçidinin metrelerce önünde durmuş oluyoruz ve görev doğru yerde
uygulanmıyor. Bunun yerine tabela algılandığında araç \ey{YAKLASMA\_SN} kadar
ilerlemeye devam ediyor, görevi sonra uyguluyor. Kameranın görüş mesafesi ve
aracın hızı sabit olduğu için bu süre tekrarlanabilir bir konum veriyor, saha
denemesiyle ayarlanacak.

Aynı tabela iki kez görülürse diye kilit koyduk. Bir tabela, araç ona
yaklaşırken onlarca ardışık karede görünüyor. Kilit olmasa araç aynı yaya
geçidinde defalarca durup kalkardı. Çözüm \ey{son\_tabela} değişkeni: bir
tabela işlendiğinde türü buraya yazılıyor, aynı tür tekrar okunduğu sürece
hiçbir şey yapılmıyor, \ey{tabela\_oku()} YOK döndürdüğünde yani tabela görüş
alanından çıktığında kilit açılıyor.

Bu yaklaşım pistte aynı türden ikinci bir tabela bulunmasını da doğru işliyor.
Araç birinci tabelayı geçiyor, arada YOK okuyor, kilit açılıyor ve ikinci tabela
normal şekilde işleniyor.

\textbf{Park yeri neye göre seçilir? Renk bilgisi hiç gelmezse araç ne yapar?}

Park yeri tek bir girdiye göre seçiliyor: sürecin başında
\ey{yuk\_rengi\_bekle()} ile aldığımız renk. PARK tabelası işlendiğinde
\ey{park\_et(renk)} çağrılıyor. Rengi bir kez alıp değişkende tuttuğumuz için
park anında yeniden okumaya veya MQTT'ye ihtiyaç yok. Bu sayede park sırasında
olası bir ağ kesintisi görevi etkilemiyor.

Renk hiç gelmezse araç zaten hiç hareket etmemiş oluyor.
\ey{yuk\_rengi\_bekle()} bloklayıcı, süreç orada duruyor ve \ey{ilerle()}
çağrısına asla ulaşılmıyor. Yani "araç yolda ama rengi bilmiyor" durumu bu
tasarımda oluşamıyor. Sorunu çalışma zamanında çözmek yerine akış yapısıyla
ortadan kaldırmış oluyoruz.

\textbf{Park tamamlanınca ne olur?}

\ey{park\_et(renk)} döndükten sonra \ey{dur()} çağırıp görev döngüsünden
çıkıyoruz. Araç süreci sona eriyor, yeni bir görev beklemiyor. Görev tanımı da
aracın tek görevlik olduğunu ve park ile birlikte duracağını söylüyor. Sondaki
\ey{dur()}, \ey{park\_et()} içindeki son manevradan sonra artık hiçbir hareket
komutunun kalmadığını garanti ediyor.

Burada iki sürecin ömrü ayrışıyor. Robot kol döngüsü devam ediyor ve konveyöre
gelen sonraki küp için yeniden başla sinyali bekliyor. Sistem bir bütün olarak
çalışmaya devam ediyor, duran yalnızca o aracın görevi.

\clearpage

\section{Güvenlik kurallarına uyum}

Görev tanımının sonunda üç kesin kural var. Aşağıdaki tablo her kuralı
algoritmada nerede ve nasıl karşıladığımızı gösteriyor.

\begin{longtable}{@{}p{4.3cm}p{3.2cm}p{7.2cm}@{}}
\toprule
\textbf{Kural} & \textbf{Karşılandığı yer} & \textbf{Nasıl} \\
\midrule
\endhead
Araç yerinde değilken yükleme yapılmamalı, başla sinyali bunu garanti ediyor. &
Robot kol, adım 1, \ey{basla\_geldi()} &
Kol araç varlığını kendisi sorgulamıyor. PLC, BP2 (küp) ve BP4 (araç)
sensörlerinin ikisini birden sağladığında sinyali veriyor. Kol yalnızca bu
sinyale bakıyor. Tek doğrulama noktası olduğu için çift kontrolden kaynaklanan
tutarsızlık riski de yok. \\
\addlinespace
MZ80 algılaması ve renk bildirimi birlikte sağlanmadan araç hareket etmemeli. &
Otonom araç, birinci ve ikinci şart &
\ey{ilerle()} çağrısına hem \ey{yuk\_rengi\_bekle()} döndükten hem de
\ey{yuk\_algilandi()} True olduktan sonra ulaşılıyor. Tek şartla ilerleyen
hiçbir yol yok. Bu bir zamanlama tercihi değil, akışın kendi yapısı. \\
\addlinespace
Kırmızı ışıkta kesinlikle geçilmemeli. &
\ey{TRAFIK\_ISIGI\_KONTROL()} &
Bekleme döngüsünden çıkış şartı yalnızca YESIL okumak. YOK okumak beklemeyi
sürdürüyor. Zaman aşımıyla geçiş yolunu bilerek koymadık. \\
\bottomrule
\end{longtable}

\section{Varsayımlar}

Aşağıdakiler görev tanımında doğrudan belirtilmeyen ama tasarımın dayandığı
noktalar. Her birinin yanında, aksi ortaya çıkarsa algoritmanın hangi tek
noktasında değişiklik gerekeceğini de yazdık.

\begin{longtable}{@{}p{0.9cm}p{6.6cm}p{7.2cm}@{}}
\toprule
\textbf{No} & \textbf{Varsayım} & \textbf{Aksi hâlde değişecek yer} \\
\midrule
\endhead
V-1 &
\ey{tabela\_oku()} dönüş kümesi YAYA, TUMSEK, PARK ve YOK olarak verilmiş,
HEMZEMIN bu kümede yok. Görev 5 kapsamında hemzemin sınıfını da eğittik ama
Görev 1'in hazır eylemi onu döndürmediği için akışta hemzemin dalı yok. &
Finalde hemzemin bir görev olarak eklenirse Şekil 3'teki tabela dallanmasına
tek bir \ey{DURUM "HEMZEMIN"} satırı ekleniyor. Kilitleme ve yaklaşma mantığı
aynen geçerli kalıyor. \\
\addlinespace
V-2 &
\ey{bekle(sn)} bloklayıcı ve bekleme sırasında araç son verilen hareket
komutunu koruyor. Yani \ey{ilerle()} sonrası \ey{bekle()} çağrısı boyunca araç
ilerlemeye devam ediyor. &
Bu yüzden yaya ve tümsek dallarında görev önce \ey{dur()} ile başlıyor. Aksi
davranışta, yani bekleme sırasında otomatik duruş olursa, yalnızca
\ey{YAKLASMA\_SN} yorumu değişiyor. \\
\addlinespace
V-3 &
Hazır eylemler arasında hız ayarı yok, yalnızca \ey{ilerle()} ve \ey{dur()}
var. Tümsek tabelasının beklenen davranışı olan hız kesmeyi, eldeki eylemlerle
en yakın karşılık olan kısa duraklamayla modelledik. &
Finalde bir hız komutu verilirse \ey{dur()}, \ey{bekle(1)}, \ey{ilerle()}
üçlüsü tek bir yavaş ilerleme çağrısıyla değiştiriliyor. \\
\addlinespace
V-4 &
\ey{basla\_geldi()} seviye sinyali. PLC'nin K10 kontağı robot çevrimi boyunca
kapalı kalıyor ve küp banttan kalkınca açılıyor. &
Sinyal darbe olsaydı çevrim sonundaki bekleme döngüsü gereksiz kalırdı. Yine de
zararsız, çünkü darbe zaten hemen düşüyor. \\
\addlinespace
V-5 &
\ey{yuk\_rengi\_bekle()} bloklayıcı ve \ey{arac/yuk} konusuna gelen ilk mesajı
döndürüyor. &
Birden fazla mesaj gelirse ilki esas alınıyor. Kol bir çevrimde tek mesaj
gönderdiği için bu durum zaten oluşmuyor. \\
\addlinespace
V-6 &
Şerit takibi aracın alt seviye kontrolünde sürekli çalışıyor, \ey{ilerle()}
"şeridi takip ederek ilerle" anlamında. Şerit takibi bu algoritmanın konusu
değil. &
Şerit kaybı ayrı bir hata durumu olarak ele alınırsa görev döngüsüne bir
kontrol adımı ekleniyor. \\
\addlinespace
V-7 &
Robot kolun GORME, AL ve YUKLE pozları önceden öğretilmiş ve aralarındaki
yörünge çarpışmasız. &
Bu üç poz hazır kabul edildiği için ara nokta tanımlamadık. \\
\bottomrule
\end{longtable}

\section{Özet}

Bu algoritmayı ayırt eden üç karar var.

Birincisi, renk bildirimini çevrimin en son adımına koymamız. Araç ancak kol
güvenli poza çekildikten sonra hareket edebiliyor, yani kol ile araç
çarpışmasını bir zamanlama ayarına değil akışın sırasına bağladık.

İkincisi, kırmızı ışıktan çıkış şartını olumlu yazmamız. "KIRMIZI değil" yerine
"YESIL" dedik, böylece algılama kaybının kural ihlaline dönüşmesi imkânsız hâle
geldi.

Üçüncüsü, tabela işlemeyi kilitli yapmamız. Bir tabela ancak görüş alanından
çıkıp yeniden girdiğinde tekrar işleniyor, aynı görev ardışık karelerde
defalarca tetiklenmiyor.

Üçünde de aynı yaklaşımı izledik: hata olasılığını çalışma anında kontrol
etmeye çalışmak yerine akışın yapısıyla ortadan kaldırmak.

\end{document}
